% !TeX root = ../main.tex
\section{Logarithmic Hooke laws and the drained response}
\label{sec:stress-reconstruction}

We prescribe the mineral and drained skeleton responses in material
logarithmic strain. Let \(\mathbf F=\mathbf R\mathbf U\) be the right
polar decomposition. Equation \eqref{eq:spherical-mineral-metric} gives
\(\bar{\mathbf U}=a^{-1/3}\mathbf U\), independently of
\(\mathbf R_A\). Taking the logarithms yields
\begin{align}
 \mathbf\varepsilon&=\log\mathbf U,
 \label{eq:skeleton-logarithmic-strain}\\
 \bar{\mathbf\varepsilon}&=\log\bar{\mathbf U},
 \label{eq:mineral-logarithmic-strain}\\
 &=\mathbf\varepsilon-\frac13\ln a\,\mathbf I.
 \label{eq:constrained-mineral-logarithmic-strain}
\end{align}
Here \(\mathbf U\) and \(\bar{\mathbf U}\) are the skeleton and
mineral right stretches. Their logarithmic traces are \(\ln J\) and
\(\ln\bar J\). We use
\(\dev\mathbf\varepsilon=\mathbf\varepsilon
-\tr(\mathbf\varepsilon)\mathbf I/3\) for the shape-changing part.
Thus \(\bar{\mathbf\varepsilon}
=\dev\mathbf\varepsilon+\ln\bar J\mathbf I/3\).

Let \(\mathbb{C}_s\) and \(\mathbb{C}^d\) be constant fourth-order
mineral and drained stiffness tensors, with the usual minor and major
symmetries and positive elastic energy. The superscript \(d\) denotes
the drained skeleton. Fourth-order tensors are set in blackboard bold;
second-order tensors are set in upright bold. The prescribed generalized
Hooke laws are
\begin{align}
 \pder{\bar W_s}{\bar{\mathbf\varepsilon}}
 &=\mathbb{C}_s:\bar{\mathbf\varepsilon},
 \label{eq:prescribed-mineral-log-hooke}\\
 \pder{W^d}{\mathbf\varepsilon}
 &=\mathbb{C}^d:\mathbf\varepsilon.
 \label{eq:prescribed-skeleton-log-hooke}
\end{align}
These stresses are conjugate to logarithmic strain. Their spatial
Kirchhoff stresses are obtained by differentiating the strain with respect
to deformation, as detailed in \cref{sec:logarithmic-derivative}.
The energies, zero in the undeformed state, are
\begin{equation}
 \bar W_s=\frac12\bar{\mathbf\varepsilon}:\mathbb{C}_s:
                    \bar{\mathbf\varepsilon},\qquad
 W^d=\frac12\mathbf\varepsilon:\mathbb{C}^d:\mathbf\varepsilon.
 \label{eq:prescribed-logarithmic-energies}
\end{equation}
The mineral energy \(\bar W_s\) is normalized per reference mineral
volume, and the drained energy \(W^d\) is normalized per reference
mixture volume. At zero pressure the coupled energy must recover \(W^d\).

We first determine the scalar distention energy from a spherical drained
strain. Define
\begin{equation}
 \mathcal{K}=\frac19\mathbf I:\mathbb{C}^d:\mathbf I,\qquad
 \mathcal{K}_s=\frac19\mathbf I:\mathbb{C}_s:\mathbf I.
 \label{eq:spherical-strain-stiffnesses}
\end{equation}
These coefficients give mean logarithmic stress per logarithmic volume
change when strain is spherical. For an anisotropic solid, a spherical
strain can generate a deviatoric stress; \(\mathcal{K}\) and \(\mathcal{K}_s\) need not
equal moduli measured under hydrostatic stress. They become the usual
bulk moduli in the isotropic limit.

Along the spherical drained path, both strains are spherical. Taking the
trace of \eqref{eq:constitutive-kirchhoff-phase-stress} at \(p=0\)
gives the same volume relation as the scalar construction of
\citet{foster2026poroplastic},
\begin{align}
 \mathcal{K}\ln J&=\phi_{s0}\mathcal{K}_s\ln\bar J,
 \label{eq:drained-mineral-volume-matching}\\
 \ln a&=\left(1-\frac{\mathcal{K}}{\phi_{s0}\mathcal{K}_s}\right)\ln J.
 \label{eq:drained-distention-volume-matching}
\end{align}
For drained spherical deformation, the mean Kirchhoff stress is
\(\tr\mathbf\tau'/3=\mathcal{K}\ln J\). The distention work in
\eqref{eq:distention-mineral-energy-work} therefore gives
\(\partial W_A/\partial\ln a=\mathcal{K}\ln J\).
Using \eqref{eq:drained-distention-volume-matching} to express \(\ln J\)
in terms of \(\ln a\), then integrating with \(W_A(1)=0\), gives
\begin{equation}
 W_A(a)=\frac{\mathcal{K}}{2\left(1-\mathcal{K}/(\phi_{s0}\mathcal{K}_s)\right)}(\ln a)^2.
 \label{eq:distention-energy}
\end{equation}
We assume \(0<\mathcal{K}<\phi_{s0}\mathcal{K}_s\), so the distention energy has positive
curvature with respect to \(\ln a\). Because \(W_A\) depends on \(a\) alone,
the energy determined from drained spherical deformation also applies when
the mineral changes shape.

Combining \eqref{eq:equivalent-mixture-mineral-energy} and
\eqref{eq:distention-energy} gives the complete elastic energy
\begin{align}
 W_s(\mathbf F,\bar J)
 &=\frac{\mathcal{K}}{2\left(1-\mathcal{K}/(\phi_{s0}\mathcal{K}_s)\right)}
             \left[\ln\left(\frac J{\bar J}\right)\right]^2
 \nonumber\\
 &\quad+\frac{\phi_{s0}}2
 \left(\dev\mathbf\varepsilon+\frac{\ln\bar J}{3}\mathbf I\right)
 :\mathbb{C}_s:
 \left(\dev\mathbf\varepsilon+\frac{\ln\bar J}{3}\mathbf I\right).
 \label{eq:equivalent-anisotropic-energy}
\end{align}
To exhibit the volume contribution, expand the mineral term:
\begin{align}
 W_s
 &=\frac{\mathcal{K}}{2\left(1-\mathcal{K}/(\phi_{s0}\mathcal{K}_s)\right)}
             \left[\ln\left(\frac J{\bar J}\right)\right]^2
 +\frac{\phi_{s0}\mathcal{K}_s}{2}(\ln\bar J)^2
 \nonumber\\
 &\quad+\frac{\phi_{s0}}3\ln\bar J\,
        \mathbf I:\mathbb{C}_s:\dev\mathbf\varepsilon
 +\frac{\phi_{s0}}2\dev\mathbf\varepsilon:\mathbb{C}_s:
                                      \dev\mathbf\varepsilon.
 \label{eq:equivalent-volumetric-energy}
\end{align}
The first two terms are precisely the companion's elastic volumetric
energy. The third describes mineral volume--shape coupling and follows
from the prescribed mineral Hooke law. It vanishes for an isotropic mineral.
The final term is the mineral shape energy, weighted by reference solid
fraction.

We enforce pressure equilibrium by differentiating this energy at fixed
\(\mathbf F\). At fixed skeleton deformation, \(J\) and the skeleton shape
remain fixed, and the distention changes only through the mineral volume.
Multiplying the derivative by \(\bar J\) gives
\begin{equation}
 -\phi_{s0}p\bar J
 =-\frac{\mathcal{K}\ln(J/\bar J)}{1-\mathcal{K}/(\phi_{s0}\mathcal{K}_s)}
   +\phi_{s0}\mathcal{K}_s\ln\bar J
   +\frac{\phi_{s0}}3\mathbf I:\mathbb{C}_s:
                                      \dev\mathbf\varepsilon.
 \label{eq:anisotropic-energy-pressure-derivative}
\end{equation}
Collecting terms yields the mineral-volume equation
\begin{equation}
 0=\mathcal{K}_s\ln\bar J
   +\left(1-\frac{\mathcal{K}}{\phi_{s0}\mathcal{K}_s}\right)p\bar J
   -\frac{\mathcal{K}}{\phi_{s0}}\ln J
   +\frac13\left(1-\frac{\mathcal{K}}{\phi_{s0}\mathcal{K}_s}\right)
      \mathbf I:\mathbb{C}_s:\dev\mathbf\varepsilon.
 \label{eq:anisotropic-mineral-eos}
\end{equation}
At fixed deformation and nonnegative pressure, the right-hand side of
\eqref{eq:anisotropic-mineral-eos} increases strictly from negative to positive
infinity as \(\bar J\) ranges over positive values. There is therefore one positive root. At negative pressure,
we follow the root continuous from zero pressure while
\begin{equation}
 \mathcal{K}_s+\left(1-\frac{\mathcal{K}}{\phi_{s0}\mathcal{K}_s}\right)p\bar J>0.
 \label{eq:trace-mineral-stability-domain}
\end{equation}
This condition keeps the local volume derivative positive. It also gives a
local minimum of \(W_s+\phi_{s0}p\bar J\) with respect to mineral volume.
At negative pressure this minimum need not be global. Positive phase volumes
and stability against general deformation must be checked separately.

Drained spherical deformation determines the distention energy \(W_A\).
Requiring the model to reproduce the entire drained Hooke law places a
further restriction on \(\mathbb{C}^d\).
Setting \(p=0\) in \eqref{eq:anisotropic-mineral-eos} gives
\begin{equation}
 \ln a=\frac1{3\mathcal{K}_s}\left(1-\frac{\mathcal{K}}{\phi_{s0}\mathcal{K}_s}\right)
                  \mathbf I:\mathbb{C}_s:\mathbf\varepsilon.
 \label{eq:drained-anisotropic-distention}
\end{equation}
To identify the drained stiffness, we first eliminate the mineral volume
from the energy. Using \(\ln\bar J=\ln J-\ln a\) in
\eqref{eq:equivalent-anisotropic-energy} replaces the mineral strain by
\(\bar{\mathbf\varepsilon}=\mathbf\varepsilon-\ln a\,\mathbf I/3\).
Expanding the mineral quadratic form gives
\begin{align}
 W_s(\mathbf F,\bar J)
 &=\frac{\phi_{s0}}2\mathbf\varepsilon:\mathbb{C}_s:\mathbf\varepsilon
 -\frac{\phi_{s0}}3\ln a\,
       \mathbf I:\mathbb{C}_s:\mathbf\varepsilon
 \nonumber\\
 &\quad+\frac12\left[
   \frac{\mathcal{K}}{1-\mathcal{K}/(\phi_{s0}\mathcal{K}_s)}
   +\phi_{s0}\mathcal{K}_s\right](\ln a)^2.
 \label{eq:drained-energy-before-elimination}
\end{align}
The two terms in square brackets come from distention and mineral volume
change, respectively. Their sum is
\(\phi_{s0}\mathcal{K}_s/[1-\mathcal{K}/(\phi_{s0}\mathcal{K}_s)]\).
Substitution of \eqref{eq:drained-anisotropic-distention} into the last
two terms therefore gives
\begin{align}
 W_s(\mathbf F,\bar J(\mathbf F,0))
 &=\frac{\phi_{s0}}2\mathbf\varepsilon:\mathbb{C}_s:\mathbf\varepsilon
 \nonumber\\
 &\quad-\frac{\phi_{s0}}{18\mathcal{K}_s}
 \left(1-\frac{\mathcal{K}}{\phi_{s0}\mathcal{K}_s}\right)
 \left(\mathbf I:\mathbb{C}_s:\mathbf\varepsilon\right)^2.
 \label{eq:drained-energy-after-elimination}
\end{align}
Here \(\bar J(\mathbf F,0)\) is the mineral volume determined by pressure
equilibrium at zero pore pressure. The negative term is the energy
reduction permitted by adjustment of mineral volume at fixed skeleton
deformation.

We require this energy to equal
\(W^d=\mathbf\varepsilon:\mathbb{C}^d:\mathbf\varepsilon/2\)
for every symmetric logarithmic strain. The major symmetry of
\(\mathbb{C}_s\) allows the squared scalar to be written as
\begin{equation}
 \left(\mathbf I:\mathbb{C}_s:\mathbf\varepsilon\right)^2
 =\mathbf\varepsilon:
 \left[(\mathbb{C}_s:\mathbf I)\otimes(\mathbb{C}_s:\mathbf I)\right]
 :\mathbf\varepsilon.
 \label{eq:drained-energy-dyadic-identity}
\end{equation}
The tensor product \(\otimes\) here forms a fourth-order tensor: its
action on \(\mathbf\varepsilon\) is \(\mathbb{C}_s:\mathbf I\)
multiplied by the scalar \(\mathbf I:\mathbb{C}_s:\mathbf\varepsilon\).
Comparing the coefficients of the quadratic forms gives
\begin{equation}
 \mathbb{C}^d=\phi_{s0}\mathbb{C}_s
 -\frac{\phi_{s0}}{9\mathcal{K}_s}\left(1-\frac{\mathcal{K}}{\phi_{s0}\mathcal{K}_s}\right)
    (\mathbb{C}_s:\mathbf I)\otimes(\mathbb{C}_s:\mathbf I).
 \label{eq:drained-stiffness-restriction}
\end{equation}
The compliance relation expresses skeleton strain as the sum of mineral
strain and the additional spherical strain from distention. To obtain the
mineral contribution, contract
\eqref{eq:drained-stiffness-restriction} with \(\mathbf\varepsilon\)
and use \eqref{eq:drained-anisotropic-distention}:
\begin{align}
 \mathbb{C}^d:\mathbf\varepsilon
 &=\phi_{s0}\mathbb{C}_s:\mathbf\varepsilon
   -\frac{\phi_{s0}}3\ln a\,\mathbb{C}_s:\mathbf I,
 \label{eq:drained-stress-distention-split}\\
 &=\phi_{s0}\mathbb{C}_s:\bar{\mathbf\varepsilon}.
 \label{eq:drained-mineral-logarithmic-stress}
\end{align}
The second line uses the strain split
\eqref{eq:constrained-mineral-logarithmic-strain}. Inverting the mineral
stiffness on symmetric tensors gives
\begin{equation}
 \bar{\mathbf\varepsilon}
 =(\phi_{s0}\mathbb{C}_s)^{-1}:\mathbb{C}^d:\mathbf\varepsilon.
 \label{eq:drained-mineral-strain-from-stress}
\end{equation}

We next express the distention in terms of the same drained logarithmic
stress \(\mathbb{C}^d:\mathbf\varepsilon\). Taking the trace of
\eqref{eq:drained-stress-distention-split} and using
\(\mathbf I:\mathbb{C}_s:\mathbf I=9\mathcal{K}_s\) gives
\begin{align}
 \mathbf I:\mathbb{C}^d:\mathbf\varepsilon
 &=\phi_{s0}\mathbf I:\mathbb{C}_s:\mathbf\varepsilon
       -3\phi_{s0}\mathcal{K}_s\ln a,
 \label{eq:drained-logarithmic-stress-trace-split}\\
 &=\frac{\mathcal{K}}{\mathcal{K}_s}
       \mathbf I:\mathbb{C}_s:\mathbf\varepsilon.
 \label{eq:drained-logarithmic-stress-trace}
\end{align}
The second line follows by substituting
\eqref{eq:drained-anisotropic-distention}. Thus
\(\mathbf I:\mathbb{C}_s:\mathbf\varepsilon\) in that equation can be
replaced by
\((\mathcal{K}_s/\mathcal{K})\mathbf I:\mathbb{C}^d:\mathbf\varepsilon\),
which gives
\begin{equation}
 \ln a=\frac{1-\mathcal{K}/(\phi_{s0}\mathcal{K}_s)}{3\mathcal{K}}
             \mathbf I:\mathbb{C}^d:\mathbf\varepsilon.
 \label{eq:drained-distention-from-stress}
\end{equation}
Substituting \eqref{eq:drained-mineral-strain-from-stress} and
\eqref{eq:drained-distention-from-stress} into
\(\mathbf\varepsilon=\bar{\mathbf\varepsilon}+\ln a\mathbf I/3\)
yields
\begin{equation}
 \mathbf\varepsilon
 =(\phi_{s0}\mathbb{C}_s)^{-1}:\mathbb{C}^d:\mathbf\varepsilon
 +\frac{1-\mathcal{K}/(\phi_{s0}\mathcal{K}_s)}{9\mathcal{K}}
       \mathbf I\left(\mathbf I:\mathbb{C}^d:\mathbf\varepsilon\right).
 \label{eq:drained-compliance-work-derivation}
\end{equation}
The first term is the mineral strain; the second is the spherical strain
added by distention. Since \(\mathbb{C}^d\) is positive definite, the
stress \(\mathbb{C}^d:\mathbf\varepsilon\) ranges over all symmetric
tensors as \(\mathbf\varepsilon\) varies. Replacing the left-hand side
by \((\mathbb{C}^d)^{-1}:\mathbb{C}^d:\mathbf\varepsilon\) and
subtracting the mineral contribution therefore identifies the compliance
difference:
\begin{equation}
 (\mathbb{C}^d)^{-1}-(\phi_{s0}\mathbb{C}_s)^{-1}
 =\frac{1-\mathcal{K}/(\phi_{s0}\mathcal{K}_s)}{9\mathcal{K}}\mathbf I\otimes\mathbf I.
 \label{eq:drained-compliance-restriction}
\end{equation}
The inverse is on symmetric tensors. Acting on a stress, the right-hand
side produces a spherical strain proportional to its trace. This is exactly
the additional deformation allowed by the dilation part of spherical distention. Together, \eqref{eq:drained-stiffness-restriction} and
\eqref{eq:drained-compliance-restriction} are necessary and sufficient within
the present construction for the full drained energy to equal the prescribed
\(W^d\). Positive mineral stiffness
and \(0<\mathcal{K}<\phi_{s0}\mathcal{K}_s\) then give positive drained stiffness. Arbitrary
independent choices of the two stiffness tensors generally do not satisfy
this relation.
